\section{Variational and Proximal Foundations of the Symmetric Integrator}

On a Hessian manifold defined by a strictly convex barrier function $\phi$, the transition from an initial state $(z_0, v_0)$ to a subsequent state $(z_1, v_1)$ over a time step $h$ can be characterized through a sequence of proximal optimizations. These updates balance the information geometry of the barrier with the mechanical momentum of the particle.

\subsection{The Primal Proximal Objective}
The primal update anchors its physics in the coordinate space $z$ and utilizes the initial Hessian $H_0 = \nabla^2 \phi(z_0)$ as the local metric. Define the \textbf{linear prediction}:
\begin{equation}
z_{\mathrm{pred}} = z_0 + 2h v_0
\end{equation}
This is the position a free particle would reach after time $2h$ at velocity $v_0$. The next position $z_1$ is the unique minimizer of:
\begin{equation}
\label{eq:primal_objective}
P(z) = D_\phi(z, z_0) + \frac{1}{2}\|z - z_{\mathrm{pred}}\|_{H_0}^2
\end{equation}
where $D_\phi(z, z_0) = \phi(z) - \phi(z_0) - \langle \nabla \phi(z_0), z - z_0 \rangle$ is the Bregman divergence.

The objective balances two competing terms:
\begin{itemize}
    \item $D_\phi(z, z_0)$: the Bregman divergence from $z_0$. This is non-negative, zero only at $z_0$, and blows up as $z$ approaches the boundary of the cone. It \textbf{guarantees that $z_1$ is interior}.
    \item $\frac{1}{2}\|z - z_{\mathrm{pred}}\|_{H_0}^2$: the squared distance to the linear prediction, measured in the frozen Hessian metric. This pulls $z_1$ toward where the geodesic should go.
\end{itemize}

\subsection{Optimality Condition}
Using the identity $\nabla_z D_\phi(z, z_0) = \nabla \phi(z) - \nabla \phi(z_0)$, setting $\nabla P(z_1) = 0$ yields:
\begin{equation}
\label{eq:primal_update}
\nabla \phi(z_1) - \nabla \phi(z_0) + H_0(z_1 - z_{\mathrm{pred}}) = 0
\end{equation}
which is equivalent to the symmetric integrator equation from \texttt{integrator.tex}:
\begin{equation}
\nabla \phi(z_1) - \nabla \phi(z_0) = H_0(2hv_0 - (z_1 - z_0))
\end{equation}

\subsection{Newton Solve}
Rearranging for numerical solution:
\begin{equation}
F(z) = \nabla \phi(z) + H_0 z - \underbrace{(\nabla \phi(z_0) + H_0 z_0 + 2h H_0 v_0)}_{\text{constant (= $2h H_0 v_0$ by log-homogeneity)}} = 0
\end{equation}
The Jacobian is $J(z) = H(z) + H_0$, which is the sum of two positive definite matrices and therefore always positive definite. This ensures Newton convergence from any feasible initial guess. Since the objective $P(z)$ is strictly convex on the cone interior and has a barrier at the boundary, the unique minimizer is always interior.

\subsection{The Dual Proximal Objective}
The dual method views the update through the momentum space $\lambda = \nabla \phi(z)$. Utilizing the conjugate barrier $\phi^*(\lambda)$ and the dual metric $H_0^{-1}$, the next dual state $\lambda_1$ minimizes:

\begin{equation}
Q(\lambda) = D_{\phi^*}(\lambda, \lambda_0) + \frac{1}{2}\|\lambda - \lambda_{\mathrm{pred}}\|_{H_0^{-1}}^2
\end{equation}
where $\lambda_{\mathrm{pred}} = \lambda_0 + 2h w_0$ with $w_0 = H_0 v_0$, and $\lambda_0 = \nabla \phi(z_0)$.

By the same reasoning, $\nabla_\lambda Q(\lambda) = 0$ yields:
\begin{equation}
\label{eq:dual_update}
z_1 - z_0 + H_0^{-1}(\lambda_1 - \lambda_0) - 2hv_0 = 0
\end{equation}

\subsection{Summary}
The symmetric integrator is a Bregman-regularized proximal step:
\begin{equation}
z_1 = \arg\min_{z \in \mathrm{int}(K)} \left\{ D_\phi(z, z_0) + \frac{1}{2}\|z - z_{\mathrm{pred}}\|_{H_0}^2 \right\}
\end{equation}
The Bregman divergence enforces interiority (no need for explicit feasibility checks). The $H_0$-quadratic term provides the geodesic-tracking force. The sum $H(z) + H_0$ in the Jacobian ensures robust Newton convergence.
